% Source-aware proof for the frozen statement
% p21p3_curve_rigidity_statement_20260826.tex.
% Date: 26 August 2026.
% This packet is not yet integrated into the book.

\paragraph{Imported facts and conventions.}
We use currents, positivity, multiplication of a positive current by a
continuous function, and the differential on currents as in
\cite[\S\S7.1--7.2, Example~7.2.12, and \S7.7]{CLD25}.  In particular,
\cite[Lemma~7.7.5]{CLD25} gives
\begin{equation}
\label{eq:p21p3-transpose-dd}
 \bigl\langle \mathrm d'\mathrm d''S,\varphi\bigr\rangle
 =
 \bigl\langle S,\mathrm d'\mathrm d''\varphi\bigr\rangle .
\end{equation}
Positive currents of bidimension $(0,0)$ are Radon measures by
\cite[Proposition~7.2.8]{CLD25}; positive currents of arbitrary
bidimension have order zero by \cite[Proposition~7.2.9]{CLD25}.  We use
smooth partitions of unity and cutoffs from
\cite[Corollaries~6.2.7--6.2.8]{CLD25}.

The local tropical argument below is the one-dimensional specialization
of the localization used in the proof of
\cite[Theorem~16.8.5]{CLD25}.  We also use
\cite[Proposition~3.3.6, Lemma~7.2.5, Lemma~12.3.6,
Theorem~12.2.4, Corollary~12.2.5,
\S14.3.2, Lemma~14.3.4, Proposition~14.4.2,
\S16.6, Proposition~16.6.3, and
\S\S18.2.5--18.2.8]{CLD25}.  Temkin reduction of germs and its interior
criterion are recalled in
\cite[\S6.13 and Theorem~6.15]{GJR26}; the proof of
\cite[Lemma~6.14]{GJR26} lifts any prescribed finite family of residue
functions to analytic germs.  The finite-fibre and calibrated
multiplicity statements for tropical skeleta are
\cite[Corollary~3.18, Remark~3.20, and Remark~4.2(7)]{GJR26}.

We do \emph{not} use \cite[Corollary~16.8.6]{CLD25}.  In particular, the
argument for the curvature measure of an arbitrary continuous
plurisubharmonic function at a type-$1$ point is supplied below.

\begin{proof}[Proof of \cref{lem:p21p3-type-one-zero-atom}]
Fix a type-$1$ point $y\in Y$.  The regular one-dimensional local ring
\(\mathcal O_{Y,y}\) has a uniformizer represented, after shrinking, by
an analytic function \(f\) with \(f(y)=0\).  Regard \(f\) as a morphism to
\(\mathbf A^{1,\mathrm{an}}_K\).  Its fibre germ at \(y\) is
zero-dimensional.  Since both \(Y\) and \(\mathbf A^{1,\mathrm{an}}_K\)
are boundaryless, the boundary-composition identity
\cite[(11.4.1.3)]{CLD25} gives
\(y\in\operatorname{Int}(Y/\mathbf A^{1,\mathrm{an}}_K)\).  Hence
\cite[Lemma~11.6.5(a)]{CLD25} makes \(f\) finite at \(y\).  It is flat
there: the induced finite map from the discrete valuation ring at \(0\)
to the regular local domain \(\mathcal O_{Y,y}\) is torsion-free, hence
free.  Now \cite[\S11.6.3]{CLD25} permits us, after shrinking source and
target, to choose a compact analytic domain \(Q_{R_0}\) containing \(y\),
an admissible radius \(R_0>0\), and a finite flat morphism of constant
degree \(d>0\)
\[
             f:Q_{R_0}\longrightarrow D(0,R_0).
\]
We choose the domains so that \(f^{-1}(0)=\{y\}\).  Finite morphisms have
empty relative boundary; the boundary-composition formula
\cite[(11.4.1.3)]{CLD25} therefore shows that the relative boundary of
\(Q_{R_0}\) in the original coordinate neighborhood lies over
\(\{|T|=R_0\}\).  Choose an admissible radius \(0<R<R_0\).
On $Q_{R_0}\setminus\{y\}$ put
\[
                         s=-\log|f|.
\]
This is a CLD-smooth tropical coordinate there and
$\mathrm d'\mathrm d''s=0$.

Fix once and for all
\[
 0<\varepsilon_0<
 \tfrac12\log(R_0/R).
\]
For admissible $0<r<R$ with \(\log(R/r)>4\varepsilon_0\), put
$s_R=-\log R$ and $s_r=-\log r$.
Choose a $C^\infty$ function
\(\rho_r:\mathbf R\to[0,1]\) which is zero on
\((-\infty,s_R-\varepsilon_0]\cup[s_r+\varepsilon_0,\infty)\), is one
on \([s_R,s_r]\), is nondecreasing on
\([s_R-\varepsilon_0,s_R]\), and is nonincreasing on
\([s_r,s_r+\varepsilon_0]\).  Define
\[
 q_r(t):=\int_{-\infty}^{t}\rho_r(\tau)\,\mathrm d\tau,
 \qquad
 H_r:=\int_{\mathbf R}\rho_r(\tau)\,\mathrm d\tau .
\]
Then
\begin{equation}
\label{eq:p21p3-radial-tent-data}
 H_r\geq \log(R/r),
 \qquad
 \int_{\mathbf R}|q''_r(t)|\,\mathrm dt=2.
\end{equation}
The function
\[
 \chi_r:=q_r(-\log|f|)
\]
on the punctured coordinate tube, extended by $0$ across the outer region
and by the constant \(H_r\) across a neighborhood of
$y$, is a genuine compactly supported CLD-smooth function.  Indeed, both
extensions are made where \(q_r\) is constant; thus no
nonsmooth maximum or piecewise affine corner remains.
Its support is contained in the fixed compact \(Q_{R_0}\), because
\(R\exp(\varepsilon_0)<R_0\).

On the transition annuli,
\begin{equation}
\label{eq:p21p3-radial-tent-dd}
 \mathrm d'\mathrm d''\chi_r
   =q''_r(s)\,
        \mathrm d' s\wedge\mathrm d'' s .
\end{equation}
On either transition band \(q''_r\) has a fixed sign, so the variation
measure of \eqref{eq:p21p3-radial-tent-dd} is obtained by replacing
\(q''_r\) with \(|q''_r|\).  Proper tropical pushforward under the finite
map \(f\), in the form of \cite[Proposition~14.4.2]{CLD25}, then counts
every source annulus with its calibrated local degree.  If
\(\Sigma_f\) is the resulting source skeleton, \(\Sigma_D\) is the radial
skeleton of the disc, and \(\mu_f,\mu_D\) are their canonical calibrated
measures, then \(f_*\mu_f=d\mu_D\).  Applying the projection formula
separately to the two positive transition bands and using
\eqref{eq:p21p3-radial-tent-data} gives the exact calculation
\begin{equation}
\label{eq:p21p3-radial-tent-tv}
 \left\|\mathrm d'\mathrm d''\chi_r\wedge\delta_Y\right\|_{\mathrm{TV}}
 =\int_{\Sigma_f}|q''_r(s)|\,\mathrm d\mu_f
 =d\int_{\Sigma_D}|q''_r(t)|\,\mathrm d\mu_D
 =d\int_{\mathbf R}|q''_r(t)|\,\mathrm dt=2d .
\end{equation}

Let \(\mu_u=\mathrm d'\mathrm d''[u]\).  It is a positive Radon measure
by \cite[Lemma~8.5.5 and Proposition~7.2.8]{CLD25}.  Positivity,
\eqref{eq:p21p3-transpose-dd}, and
\eqref{eq:p21p3-radial-tent-tv} yield
\begin{align}
 \log(R/r)\,\mu_u(\{y\})
 &\leq
    \int_Y\chi_r\,\mathrm d\mu_u                                                       \notag\\
 &= \bigl\langle [u],
            \mathrm d'\mathrm d''\chi_r\bigr\rangle                                    \notag\\
 &\leq
   \left|\bigl\langle [u],
            \mathrm d'\mathrm d''\chi_r\bigr\rangle\right|                             \notag\\
 &\leq
   2d\sup_{Q_{R_0}}|u|.
\label{eq:p21p3-radial-tent-bound}
\end{align}
The last quantity is independent of \(r\).  The nontriviality of the
valuation supplies admissible radii tending to zero.  Letting \(r\) tend to zero
in \eqref{eq:p21p3-radial-tent-bound} gives
\(\mu_u(\{y\})=0\).
\end{proof}

\begin{proof}[Proof of \cref{prop:p21p3-curve-rigidity}]
We first record the elementary distribution calculation that will be
used in every tropical chart.  Let \(\Gamma\) be a locally finite
one-dimensional polyhedral set in an open subset of \(\mathbf R^N\), and
let \(S\) be a positive current of bidimension \((1,1)\), supported on
\(\Gamma\), with \(\mathrm d' S=\mathrm d''S=0\).  In Euclidean
coordinates, the order-zero assertion of
\cite[Proposition~7.2.9]{CLD25} and positivity write \(S\) as a
symmetric positive
semidefinite matrix \(M=(M_{jk})\) of Radon measures; closedness is the
distributional system
\begin{equation}
\label{eq:p21p3-matrix-divergence}
                   \sum_j\partial_jM_{jk}=0
                   \quad (1\leq k\leq N).
\end{equation}
Consider the relative interior of an edge and use affine coordinates
\((t,z_2,\ldots,z_N)\) in which the edge is
\(\{z_2=\cdots=z_N=0\}\).  Testing
\eqref{eq:p21p3-matrix-divergence} against
\(\psi(t)z_\ell\), with \(\ell>1\), gives \(M_{\ell k}=0\) on that
edge.  Symmetry then leaves only \(M_{11}\), and testing against
\(\psi(t)\) gives \(\partial_tM_{11}=0\).  Thus
\begin{equation}
\label{eq:p21p3-edge-density}
                     S|_e=b_e[\,e\,],
                     \qquad b_e\geq0,
\end{equation}
where \(b_e\) is constant on the edge.  Testing
\eqref{eq:p21p3-matrix-divergence} with a function constant near a vertex
isolates the endpoint terms of the adjacent edge measures and gives the
vector balancing relation
\begin{equation}
\label{eq:p21p3-current-balancing}
                       \sum_{e\ni v}b_eu_{e,v}=0.
\end{equation}
Here \(u_{e,v}\) is the primitive outgoing tangent.  A possible positive
remainder at a vertex has the form \(P\delta_v\) for a positive
semidefinite matrix \(P\).
After the preceding zero-order endpoint terms have balanced, substitution
in \eqref{eq:p21p3-matrix-divergence} and testing against functions that
vanish at \(v\) but have arbitrary first jet there isolates the
derivatives of \(P\delta_v\) and gives \(P=0\).
This proves the calculation,
including for a finite substar of an infinite-valence star.

\smallskip
\emph{The rig-smooth case.}
Assume first that \(X\) is connected and rig-smooth.  We spell out the
localization that prevents a tropical pushforward from merging different
source branches.  Fix a compactly supported smooth test form.  A finite
partition of unity and \cite[Lemma~12.3.6]{CLD25} reduce it to finitely
many summands with the following properties: a summand is tropicalized by
a moment \(F_0\) on a compact analytic domain \(V\), and its tropical
support is disjoint from \(F_{0,\mathrm{trop}}(\partial V)\).  If necessary,
concatenate \(F_0\) with the moment which tropicalizes the summand.  Over
the relative interior of a one-cell, the relevant tropical-skeleton fibre
is finite by
\cite[Corollary~3.18 and Remark~3.20]{GJR26}; after subdividing, only
finitely many source one-cells and vertex stars meet the support, because
the subdivision is locally finite and the support is compact.

Localize once more in the source, before adding any analytic germs.  On
the relative interiors of the finitely many source one-cells choose
mutually disjoint smooth plateaus \(\lambda_j\), and decompose the summand
accordingly; concatenating \(F_0\) with moments defining the
\(\lambda_j\) tropicalizes the new summands.  Near a relevant vertex keep
a finite source star instead of deleting the vertex.  At a type-$2$
vertex \(x\), Riemann--Roch supplies rational functions on its smooth
residual curve whose order vectors distinguish any prescribed finite set
of directions.  More precisely, for selected closed points
\(p_1,\ldots,p_s\), choose an effective divisor \(D\) away from them with
\(H^1(C_x,\mathcal O_{C_x}(D-2\sum_jp_j))=0\).  The divisor exact sequence
makes the evaluation-and-first-jet map from
\(H^0(C_x,\mathcal O_{C_x}(D))\) surjective.  We can thus choose, for each
selected point, a rational function with a simple zero there and a
nonzero value at all the other selected points.  On the already localized
chart these functions lift to analytic units by the construction in the
proof of \cite[Lemma~6.14]{GJR26}; append the lifts to the local moment.
At a type-$3$ vertex append an annulus coordinate.  A finite repetition
and a subdivision make the refined moment injective on every relevant
source one-cell and distinguish all relevant germs at every retained
vertex.  Appending coordinates does not alter the form: the old moment is
the projection of the refined one.  Notice that the residue lifts are
never required to extend across the original global domain; source
localization precedes their construction.

Fix one of these finitely many localized summands, write \(F\) for its
refined moment, and retain the notation \(V\) for its compact source
domain.  Put
\[
 \Omega=\mathbf R^N\setminus F_{\mathrm{trop}}(\partial V),
 \qquad U=F_{\mathrm{trop}}^{-1}(\Omega).
\]
Exactly as in the base-change argument in the proof of
\cite[Theorem~16.8.5]{CLD25}, the boundary-avoidance construction makes
\(F_{\mathrm{trop}}|_U:U\to\Omega\) proper, and
\[
 \langle T_\Omega,\beta\rangle
       :=\langle T|_U,F^*\beta\rangle,
 \qquad \beta\in A_c^{1,1}(\Omega),
\]
is a positive closed Euclidean current supported on a locally finite
one-dimensional polyhedral set.  On every source one-cell the calculation
\eqref{eq:p21p3-edge-density} gives a constant nonnegative flux \(b_e\).
If \(m_e>0\) is the flux of the canonical integration current on that same
source one-cell, define
\begin{equation}
\label{eq:p21p3-branch-ratio}
                              a_e:=b_e/m_e.
\end{equation}
On a common injective refinement, the elementary one-dimensional
pushforward formula for the matrix measure in
\eqref{eq:p21p3-edge-density} multiplies the ordinary-current flux \(b_e\)
by the expansion index of the tangent lattice.  The canonical flux \(m_e\)
is multiplied by the same index, by
\cite[\S14.3.2, Lemma~14.3.4, and Proposition~14.4.2]{CLD25}.  Explicitly,
\[
 F_{\mathrm{trop},*}v_e=q_eu_e
 \quad\Longrightarrow\quad
 b_e=q_e\bar b_e,\qquad m_e=q_e\bar m_e,
 \qquad \frac{b_e}{m_e}=\frac{\bar b_e}{\bar m_e},
\]
where barred coefficients are measured in the source primitive
coordinate.  Hence
\(a_e\) is independent of the moment, the subdivision, and the localized
presentation.  It is
constant along the annular collar represented by \(e\).  If a localized
summand meets no vertex, the edge calculation determines its pairing.  If
it meets a vertex, the edge densities, the balancing relation
\eqref{eq:p21p3-current-balancing}, and the equality \(P\delta_v=0\) proved
above determine the pairing on the entire finite star.  The finite
localized decomposition therefore shows that these intrinsic branch
ratios determine \(T\) on every compactly supported test form; no hidden
sheet has been summed into another one.

At a type-$3$ point the two annular directions have opposite calibrated
tangents.  Applying \eqref{eq:p21p3-current-balancing} to \(T\) and to
\(\delta_X\) shows that their ratios \(a_e\) agree.  It remains to
compare the directions at a type-$2$ point \(x\).

Apply Temkin's interior criterion
\cite[Theorem~6.15]{GJR26} to the structure morphism
\(X\to\mathcal M(K)\).  The reduction of \(\mathcal M(K)\) is its
unique point, whose inverse image in the Riemann--Zariski space of
\(\widetilde{\mathcal H(x)}\) is the whole space.  Since
\(x\in\operatorname{Int}(X)\), the criterion therefore gives
\begin{equation}
\label{eq:p21p3-full-germ-reduction}
 \operatorname{red}(X,x)
   =\mathbf P_{\widetilde{\mathcal H(x)}/\widetilde K}.
\end{equation}
For a type-$2$ point the ordinary residue field
\(\widetilde{\mathcal H(x)}\) has transcendence degree one over the
algebraically closed residue field \(\widetilde K\).  Hence it is the
function field of a unique smooth projective integral curve \(C_x\).
Under the identification
\(
 \mathbf P_{\widetilde{\mathcal H(x)}/\widetilde K}\simeq C_x
\), the trivial valuation is the generic point, while every nontrivial
valuation ring is the discrete valuation ring
\(\mathcal O_{C_x,p}\) of a unique closed point \(p\).  The
reduction-of-germs/domain correspondence
\cite[\S11.8.6]{CLD25}, together with the annular description in
\cite[\S\S18.2.5--18.2.8]{CLD25}, identifies that closed point with the
corresponding tangent direction of the punctured germ at \(x\).  Thus
\eqref{eq:p21p3-full-germ-reduction} says exactly that no closed point,
and hence no tangent direction, is missing.

We normalize the flux in each direction before using it in a divisor.
For a closed point \(p\in |C_x|_{\mathrm{cl}}\), the Riemann--Roch
jet-surjectivity argument used above supplies a residue function \(s_p\)
with a simple zero at \(p\).  Lift \(s_p\) to an analytic unit near \(x\)
and retain it as a coordinate of the local moment.  The image tangent of
the \(p\)-direction then has an integral coordinate equal to one, hence is
primitive.  Let \(m_p>0\) be the canonical calibrated flux of
\(\delta_X\) in such a normalized moment, and let \(a_p\) be the quotient
\eqref{eq:p21p3-branch-ratio}.  These numbers do not depend on the
normalized moment: on a common refinement retaining the simple-zero
coordinate, the tangent expansion index remains one, and the pushforward
comparison preceding the type-$3$ argument applies.

Define additive maps on the divisor group of \(C_x\) by
\begin{equation}
\label{eq:p21p3-divisor-functionals}
 L_T([p])=a_pm_p,\qquad L_X([p])=m_p.
\end{equation}
Both maps are nonnegative on effective divisors.

We claim that both maps vanish on principal divisors, and include the
pointwise slope calculation.  Let
\(g\in\widetilde{\mathcal H(x)}^\times
      =\operatorname{Rat}(C_x)^\times\).
Choose a norm-one representative of \(g\) in the valuation ring of
\(\mathcal H(x)\).  The
lifting step in the proof of \cite[Lemma~6.14]{GJR26} gives an analytic
germ \(f\in\mathcal O_{X,x}\) having residue \(g\).  Its value at \(x\)
has norm one, so, after shrinking, \(f\) is an analytic unit.

Fix \(p\in |C_x|_{\mathrm{cl}}\).  Choose an adapted annulus parameter \(t\),
increasing toward \(x\), for the corresponding direction as in
\cite[\S18.2.8]{CLD25}; its reduction
\(\tau\) is a uniformizer at \(p\).  In the fraction field of the
discrete valuation ring \(\mathcal O_{C_x,p}\), write
\(
 g=\tau^{r}u
\), where \(r=\operatorname{ord}_p(g)\) and \(u\) is a unit.  On the
punctured collar, \(t\) is invertible and the analytic unit
\(h:=ft^{-r}\) has reduction \(u\).  Since
\(u,u^{-1}\in\mathcal O_{C_x,p}\), the point \(p\) lies in the
quasi-compact reduction open cut out by \(u\) and \(u^{-1}\).  By the
germ-domain dictionary of \cite[\S11.8.6]{CLD25}, this open corresponds,
after shortening the collar, to
\[
                   \{|h|\leq1,\ |h^{-1}|\leq1\}
                   =\{|h|=1\}.
\]
Hence \(-\log|f|=r(-\log|t|)\) on that collar.  Since
\(-\log|t|\) has outgoing slope one, we obtain
\begin{equation}
\label{eq:p21p3-residue-slope}
 \operatorname{slope}_{p}(-\log|f|)
      =\operatorname{ord}_p(g).
\end{equation}
This also proves that only the finitely many directions in
\(\operatorname{div}(g)\) contribute.

Reapply the source-separation construction to the finite set
\(\operatorname{Supp}(\operatorname{div}(g))\), retaining for every
contributing \(p\) a lifted separator \(s_p\) with a simple zero at \(p\),
and append the analytic unit \(f\) as one further coordinate.  In the
\(p\)-direction the \(s_p\)-coordinate of the full tangent vector is one;
the image tangent is therefore primitive and its expansion index is
exactly one.  Take the scalar
\(-\log|f|\)-coordinate of the vector balancing relation
\eqref{eq:p21p3-current-balancing}.  The coefficient of the direction
corresponding to \(p\) is its slope
\eqref{eq:p21p3-residue-slope}, while every direction outside the finite
support of \(\operatorname{div}(g)\) has scalar slope zero.  Since the
current flux in that direction is \(b_p=a_pm_p\), this scalar component is
\begin{equation}
\label{eq:p21p3-T-principal}
       0=\sum_p a_pm_p\operatorname{ord}_p(g)
        =L_T(\operatorname{div}(g)).
\end{equation}
The same scalar component with \(T=\delta_X\), which is closed by
\cite[\S7.7.4]{CLD25}, gives
\begin{equation}
\label{eq:p21p3-X-principal}
       0=\sum_p m_p\operatorname{ord}_p(g)
        =L_X(\operatorname{div}(g)).
\end{equation}
Thus no cutoff on the possibly infinite-valence star is used: only the
finite set of nonzero scalar slopes enters the balancing identity.

We now use Riemann--Roch on \(C_x\).  Put
\[
                  g_x:=\dim H^1(C_x,\mathcal O_{C_x}).
\]
Let
\[
                         L:\operatorname{Div}(C_x)\longrightarrow\mathbf R
\]
be additive, nonnegative on effective divisors, and zero on principal
divisors.  If \(D\) has degree zero, choose an effective divisor \(A\)
with
\(\deg A+1-g_x>0\).  For every \(n\in\mathbf Z\), Riemann--Roch gives
\[
 \begin{split}
 h^0(C_x,\mathcal O_{C_x}(A+nD))
 &=\deg A+1-g_x
   +h^0(C_x,\omega_{C_x}(-A-nD))\\
 &\geq \deg A+1-g_x>0.
 \end{split}
\]
Thus \(A+nD\) is linearly equivalent to an effective divisor and
\[
                         0\leq L(A)+nL(D)
                         \qquad(n\in\mathbf Z).
\]
Letting \(n\) tend to both signs of infinity gives \(L(D)=0\).  Hence
\(L\) factors through degree.  Since \(\widetilde K\) is algebraically
closed, \(C_x\) has closed points of degree one, and therefore
\begin{equation}
\label{eq:p21p3-degree-classification}
                         L=c_L\deg
\end{equation}
for some \(c_L\geq0\).  Applying this first to \(L_X\) and then to
\(L_T\) yields
\[
                  m_p=c_X\deg(p),\qquad
                  a_pm_p=c_T\deg(p).
\]
Here \(c_X>0\), since every \(m_p\) is positive.  Consequently
\begin{equation}
\label{eq:p21p3-equal-star-ratios}
                              a_p=c_T/c_X
\end{equation}
for every closed point \(p\).  This proves equality of all branch
constants at every
type-$2$ point.  The argument uses only finitely many directions at a
time and is therefore valid at an infinite-valence star.

The type-$2$ and type-$3$ comparisons propagate the ratio through the
non-endpoint core of the connected rig-smooth curve.  This core is
connected.  Indeed, connected analytic spaces are arcwise connected as
recalled in \cite[\S12.1.12]{CLD25}, and the local triangulated-curve
description in \cite[\S\S18.2.5--18.2.8]{CLD25} shows that a type-$1$ or
type-$4$ endpoint cannot be an interior point of an arc joining two
non-endpoints.  Let \(E\) be the
endpoint locus, consisting of the type-$1$ and type-$4$ points.  Such a
point has tropical dimension zero.  Hence
\cite[Proposition~3.3.6]{CLD25} says, pointwise and at the level of
germs, that every smooth test form \(\omega\) of bidegree \((1,1)\)
vanishes on a neighborhood of every point of \(E\); equivalently,
\[
                         \operatorname{supp}(\omega)\cap E=\varnothing .
\]
After subtracting the constant multiple of \(\delta_X\) already found
on the non-endpoint core, the remainder current is supported in \(E\).
By the definition of the support of a current, its pairing with every
such \(\omega\) is therefore zero.  This rules out not only atoms but
also any diffuse remainder on an uncountable type-$4$ endpoint locus.
We have proved
\begin{equation}
\label{eq:p21p3-smooth-classification}
                              T=c\delta_X
\end{equation}
for one \(c\geq0\).  No compactness of \(X\) was used: every current test
has compact support, and all preceding arguments are local on that
support.

\smallskip
\emph{The reduced case.}
Return to the reduced curve in the statement.  Locally write
\[
 \mathcal Q:=\nu_*\mathcal O_{\widetilde X}/\mathcal O_X,
 \qquad
 \mathfrak c:=\operatorname{Ann}_{\mathcal O_X}(\mathcal Q).
\]
Let \(Z\) be the union of
\(\operatorname{Supp}(\mathcal Q)=V(\mathfrak c)
=\operatorname{Supp}(\mathcal O_X/\mathfrak c)\) and the pairwise
intersections of distinct irreducible components.  Here are the algebraic
details.  On a good affinoid domain the analytic local rings are excellent
by \cite[\S11.1.10]{CLD25}; normalization is finite and its components
correspond to the irreducible components by
\cite[\S11.3.1]{CLD25}.  The normalization quotient and its annihilator,
the conductor, are therefore coherent.  The quotient vanishes, and the
conductor is the unit ideal, at every generic point of the reduced curve;
thus the quotient support, or equivalently the conductor zero locus, is
zero-dimensional.  The
same dimension argument applies to the proper intersections of distinct
one-dimensional components.  Since irreducible components are locally
finite by \cite[\S11.2.5]{CLD25}, these affinoid descriptions glue and
show that \(Z\) is closed and locally finite.  Strictness over the
algebraically closed field \(K\) implies that every point of \(Z\), and
every point of
\(
 Z_i:=\nu_i^{-1}(Z)
\), is of type~$1$.  The normalization \(\nu_i\) is finite, hence
topologically proper and without relative boundary, and it is an
isomorphism away from \(Z_i\).  The closed immersion \(X_i\hookrightarrow
X\) is finite as well.  Applying the boundary-composition identity
\cite[(11.4.1.3)]{CLD25} first to \(X_i\to X\to\mathcal M(K)\) and then to
\(\widetilde X_i\to X_i\to\mathcal M(K)\) proves that
\(\widetilde X_i\) is boundaryless.  It is connected because components
of the normalization correspond to irreducible components of \(X\), as
just cited.

The fundamental-cycle identity is
\begin{equation}
\label{eq:p21p3-fundamental-cycle}
             \delta_X=\sum_i(\nu_i)_*\delta_{\widetilde X_i}.
\end{equation}
Indeed, since \(X\) is reduced, its fundamental cycle is the locally
finite sum of its irreducible components with multiplicity one.
Integration agrees with the associated cycle by
\cite[Proposition~16.6.3]{CLD25}, and normalization has calibrated
degree one over the normal locus.  This proves
\eqref{eq:p21p3-fundamental-cycle} by testing on that locus; its
complement is the type-$1$ set \(Z\), on which every $(1,1)$-test-form
germ vanishes by \cite[Proposition~3.3.6]{CLD25}.

On \(X\setminus Z\), the rig-smooth case gives
\[
                         T=\sum_i c_i\delta_{X_i},
                         \qquad c_i\geq0.
\]
There is a single coefficient on each component: deleting the locally
finite type-$1$ endpoint set \(Z_i\) does not disconnect
\(\widetilde X_i\).  Indeed, two remaining points can be joined by an arc
by \cite[\S12.1.12]{CLD25}, and the curve structure of
\cite[\S\S18.2.5--18.2.8]{CLD25} shows that a type-$1$ endpoint cannot be
an interior point of that arc.
Set
\(
 S:=T-\sum_i c_i\delta_{X_i}
\).  Then \(\operatorname{supp}S\subseteq Z\).  If
\(\alpha\in A_c^{1,1}(X)\), Proposition~3.3.6 of \cite{CLD25} makes
the germ of \(\alpha\) zero at every point of \(Z\).  By the definition
of the support of a smooth form, this says directly that
\(\operatorname{supp}\alpha\cap Z=\varnothing\), and
\(\langle S,\alpha\rangle=0\).  Thus \(S=0\), proving
\eqref{eq:p21p3-curve-classification}.  Uniqueness follows by
restriction to an open subset of \(X_i\) meeting no other component.

\smallskip
\emph{The Bedford--Taylor assertion.}
Put \(\widetilde W_i=\nu_i^{-1}(W)\).  For every continuous function
\(F\) on \(W\), the test-current identity corresponding to
\eqref{eq:p21p3-fundamental-cycle} is
\begin{equation}
\label{eq:p21p3-function-cycle-decomposition}
 [F]=F\delta_W
     =\sum_i(\nu_i)_*
        \bigl((\nu_i^*F)\delta_{\widetilde W_i}\bigr).
\end{equation}
To verify it, pair both sides with
\(\alpha\in A_c^{1,1}(W)\).  All integration currents in
\eqref{eq:p21p3-fundamental-cycle} are positive and hence have order zero.
We may therefore multiply that current identity by the continuous
function \(F\), using \cite[Example~7.2.12]{CLD25}.  For a finite map
\(\nu_i\), the identity
\[
 F(\nu_i)_*S=(\nu_i)_*((F\circ\nu_i)S)
\]
follows immediately by testing against the smooth form \(\alpha\); only
the continuous coefficient is multiplied into the order-zero measure.
This proves \eqref{eq:p21p3-function-cycle-decomposition}.  The sum is
finite on \(\operatorname{supp}\alpha\).

Pullback preserves continuous plurisubharmonic functions by
\cite[Proposition~8.1.5]{CLD25}.  On \(\widetilde W_i\), set
\[
 \mu_i^U=\mathrm d'\mathrm d''[\nu_i^*U],
 \qquad
 \mu_i^V=\mathrm d'\mathrm d''[\nu_i^*V].
\]
These are positive Radon measures by
\cite[Lemma~8.5.5]{CLD25}.  Proper pushforward commutes with the
differential on currents by transposition; equivalently, apply the
projection formula \cite[Proposition~8.5.6]{CLD25} to
\(S=\delta_{\widetilde W_i}\).  Thus
\eqref{eq:p21p3-function-cycle-decomposition}, applied first to \(U\)
and then to \(V\), transforms
\eqref{eq:p21p3-curve-equal-curvature} into
\begin{equation}
\label{eq:p21p3-component-measure-sum}
                    \sum_i(\nu_i)_*(\mu_i^U-\mu_i^V)=0.
\end{equation}
Away from \(Z\), the component images are disjoint and \(\nu_i\) is an
analytic isomorphism.  Therefore each signed measure
\(\mu_i^U-\mu_i^V\) is supported on the locally finite type-$1$ set
\(Z_i\cap\widetilde W_i\).  By
\cref{lem:p21p3-type-one-zero-atom}, neither \(\mu_i^U\) nor
\(\mu_i^V\) has an atom there.  A signed Radon measure supported on a
locally finite set is the locally finite sum of its point masses, so it
is zero if every point mass is zero.  Hence
\begin{equation}
\label{eq:p21p3-component-measure-equality}
                              \mu_i^U=\mu_i^V
                              \quad\text{for every }i.
\end{equation}
Finally put \(h=U-V\).  Multiplication \(hT\) is defined by
\cite[Example~7.2.12]{CLD25}, and
\eqref{eq:p21p3-curve-classification} gives the test-current identity
\[
 hT|_W=\sum_i c_i(\nu_i)_*
       \bigl((\nu_i^*h)\delta_{\widetilde W_i}\bigr).
\]
Proper pushforward commutes with \(\mathrm d'\mathrm d''\), so
\eqref{eq:p21p3-component-measure-equality} gives
\begin{align*}
 \mathrm d'\mathrm d''(hT|_W)
 &=\sum_i c_i(\nu_i)_*
       \mathrm d'\mathrm d''[\nu_i^*U-\nu_i^*V] \\
 &=\sum_i c_i(\nu_i)_*(\mu_i^U-\mu_i^V)=0.
\end{align*}
This is \eqref{eq:p21p3-curve-bt-conclusion}, and the last assertion is
exactly \cite[Definition~8.6.2]{CLD25}.
\end{proof}
